\documentclass[a4paper,11pt]{article}

%%% Packages
\usepackage[utf8]{inputenc}
\usepackage[T1]{fontenc}
\usepackage{amsmath,amssymb,amsthm}
\usepackage{mathtools}
\usepackage{stmaryrd}       % \llbracket, \rrbracket
\usepackage{booktabs}
\usepackage{array}
\usepackage{enumitem}
\usepackage{hyperref}
\usepackage[margin=2.5cm]{geometry}
\usepackage{xcolor}

%%% Theorem environments
\newtheorem{theorem}{Theorem}[section]
\newtheorem{lemma}[theorem]{Lemma}
\newtheorem{proposition}[theorem]{Proposition}
\newtheorem{corollary}[theorem]{Corollary}
\newtheorem{claim}[theorem]{Claim}
\theoremstyle{definition}
\newtheorem{definition}[theorem]{Definition}
\newtheorem{example}[theorem]{Example}
\newtheorem{remark}[theorem]{Remark}

%%% Notation macros (matching companion paper)
\newcommand{\ALC}{\ensuremath{\mathcal{ALC}}}
\newcommand{\ALCI}{\ensuremath{\mathcal{ALCI}}}
\newcommand{\ALCIRCC}[1]{\ensuremath{\mathcal{ALCI}_{\mathit{RCC#1}}}}
\newcommand{\NR}{\ensuremath{N_R}}
\newcommand{\NC}{\ensuremath{N_C}}
\newcommand{\DeltaI}{\ensuremath{\Delta^{\mathcal{I}}}}
\newcommand{\interp}{\ensuremath{\mathcal{I}}}
\newcommand{\DR}{\ensuremath{\mathsf{DR}}}
\newcommand{\PO}{\ensuremath{\mathsf{PO}}}
\newcommand{\EQ}{\ensuremath{\mathsf{EQ}}}
\newcommand{\PP}{\ensuremath{\mathsf{PP}}}
\newcommand{\PPI}{\ensuremath{\mathsf{PPI}}}
\newcommand{\inv}{\ensuremath{\mathsf{inv}}}
\newcommand{\comp}{\ensuremath{\mathsf{comp}}}
\newcommand{\cl}{\ensuremath{\mathsf{cl}}}
\newcommand{\sub}{\ensuremath{\mathsf{sub}}}
\newcommand{\Tp}{\ensuremath{\mathsf{Tp}}}
\newcommand{\DN}{\ensuremath{\mathsf{DN}}}
\newcommand{\EXPTIME}{\textsc{ExpTime}}
\newcommand{\PSPACE}{\textsc{PSpace}}
\newcommand{\NP}{\textsc{NP}}


\title{Closing the Extension Gap:\\
  Decidability of $\ALCIRCC{5}$ via Direct Model Construction}

\author{Claude\thanks{Claude is an AI assistant by Anthropic.
  This research was prompted by Michael Wessel (\texttt{miacwess@gmail.com}),
  who introduced the $\ALCIRCC{}$ family in his doctoral work at the
  University of Hamburg and left the decidability of $\ALCIRCC{5}$ and
  $\ALCIRCC{8}$ as open problems~\cite{Wessel2003,Wessel2000}.
  This paper is a companion to~\cite{CompanionPaper}, which introduced
  the quasimodel method and tableau calculus but left the
  soundness direction (open branch $\Rightarrow$ model) incomplete.}}

\date{April 2026}

\usepackage[normalem]{ulem}    % \sout for strikethrough

\begin{document}
\maketitle

\begin{center}
\fbox{\parbox{0.92\textwidth}{%
\textbf{Disclaimer.}\; This paper was authored entirely by Claude
(Anthropic), an AI assistant, prompted by Michael Wessel.  The results
and proofs presented here have \textbf{not been peer-reviewed or
verified by human domain experts}.  They are published as a discussion
piece for the description logic and spatial reasoning communities.
The claims should not be taken as established results unless
independently verified.}}
\end{center}

\medskip

\begin{center}
\fbox{\parbox{0.92\textwidth}{%
\textbf{RETRACTED (April 2026).}\;
This paper contains \textbf{two errors} identified by GPT-5.4 Pro's
technical review:

\smallskip
\textbf{Error 1: Algebraic error in Lemma~\ref{lem:unique-failure}.}\;
The paper states $\PP \notin \comp(\DR, \PP) = \{\DR\}$, but the
correct value is $\comp(\DR, \PP) = \{\DR, \PO, \PP\}$ (and
$\PP \in \comp(\DR, \PP)$).  The paper confused $\comp(\DR, \PP)$
with $\comp(\PP, \DR) = \{\DR\}$.  There is only \textbf{one}
self-absorption failure pair involving $\PP$: $(\PP, \DR)$, not
$(\DR, \PP)$.  The dual stated at line~305 is wrong; the correct
dual is $\PP \notin \comp(\PP, \DR) = \{\DR\}$.

\smallskip
\textbf{Error 2: Theorem~\ref{thm:path-consistent}
(path-consistency) is false.}\;
The $\DN_{\mathsf{safe}}$ domains aggregate relations from
\emph{different representatives} of the same type.  When two nodes
of type $\tau_1$ have different relations ($\PO$ and $\DR$) to a
node of type $\tau_2$, $\DN_{\mathsf{safe}}(\tau_1, \tau_2) =
\{\PO, \DR\}$.  Path-consistency fails: for a third type $\tau_3$
with $\DN_{\mathsf{safe}}(\tau_2, \tau_3) = \{\PP\}$,
$\comp(\PO, \PP) = \{\PO, \PP\} \ni \PP$ but
$\comp(\DR, \PP) = \{\DR, \PO, \PP\}$, and if
$\DN_{\mathsf{safe}}(\tau_1, \tau_3) = \{\PP\}$ then
$\{\PP\} \cap \comp(\DR, \PP) = \{\PP\} \neq \emptyset$ but
$\{\PP\} \cap \comp(\DR, \PO) = \emptyset$---the network can fail
arc-consistency.  GPT-5.4 provided a concrete counterexample
confirming this.

\smallskip
\emph{The decidability claim (Theorem~\ref{thm:main}) is
\textbf{wrong}.  The main theorem, the complexity corollary, and
the soundness theorem are all retracted.}

\smallskip
\emph{What survives.}\;
The root cause analysis (Section~\ref{sec:root-cause}) is
correct after fixing the algebraic error: there is ONE
self-absorption failure $\PPI \notin \comp(\DR, \PPI) = \{\DR\}$
(and its dual $\PP \notin \comp(\PP, \DR) = \{\DR\}$).
Lemma~\ref{lem:forced-DR} (forced $\DR$ edge) and
Theorem~\ref{thm:self-safe} (witness relations are self-safe) are
correct.  These are useful building blocks for future work.}}
\end{center}

\bigskip

%% ============================================================
\begin{abstract}
The companion paper~\cite{CompanionPaper} presented two decision
procedures for concept satisfiability in $\ALCIRCC{5}$---a
type-elimination algorithm and a tableau calculus with blocking---but
identified an \emph{extension gap}: the soundness direction of the
tableau (open branch $\Rightarrow$ model) relied on a Henkin
construction whose correctness for abstract quasimodels was unproven.
We attempted to close this gap by identifying the algebraic root
cause---the self-absorption failure
$\PPI \notin \comp(\DR, \PPI) = \{\DR\}$---and constructing a model
directly from open completion graphs via the \emph{full RCC5
tractability} theorem.
\textbf{Erratum:} This paper is \textbf{retracted}.  The algebraic
analysis contains an error ($\comp(\DR, \PP) \neq \{\DR\}$), and the
path-consistency theorem (Theorem~\ref{thm:path-consistent}) is
\textbf{false}: the $\DN_{\mathsf{safe}}$ domains are too coarse.
See the status box for details.
\end{abstract}

\medskip
\noindent\textbf{Keywords:} description logic, spatial reasoning,
Region Connection Calculus, decidability, tableau, patchwork property,
self-absorption

%% ============================================================
\section{Introduction}
\label{sec:intro}

The description logic $\ALCIRCC{5}$ extends $\ALCI$ with role names
drawn from the five RCC5 base relations $\{\DR, \PO, \EQ, \PP, \PPI\}$,
interpreted over complete-graph models satisfying the RCC5 composition
table~\cite{Wessel2003,Wessel2000}.  Decidability of concept
satisfiability in $\ALCIRCC{5}$ was left open by
Wessel~\cite{Wessel2003}, who established \PSPACE-hardness.

The companion paper~\cite{CompanionPaper} introduced two candidate
decision procedures---a type-elimination algorithm and a tableau
calculus with equality anywhere-blocking---and proved termination and
completeness (satisfiable $\Rightarrow$ open branch) for both.
However, the \emph{soundness} direction (open branch $\Rightarrow$
satisfiable) contained a gap: it proceeded by extracting a quasimodel
from the open completion graph and invoking a Henkin construction to
build an actual model.  The Henkin construction at each step solves a
constraint satisfaction problem (CSP) over a disjunctive RCC5 network,
and while the \emph{patchwork property}~\cite{RenzNebel1999} and
\emph{full RCC5 tractability}~\cite{Renz1999} guarantee solvability
for path-consistent networks, the extension step may destroy
path-consistency.  This is the \emph{extension gap}.

\subsection{Contribution}

We resolve the extension gap by:
\begin{enumerate}[nosep]
  \item \textbf{Root cause analysis} (Section~\ref{sec:root-cause}).
    We identify the unique algebraic rigidity in the RCC5 composition
    table responsible for the gap: the \emph{self-absorption failure}
    $\PPI \notin \comp(\DR, \PPI) = \{\DR\}$ (and its dual
    $\PP \notin \comp(\PP, \DR) = \{\DR\}$).  Geometrically, this
    expresses the \emph{containment collapse}: if region $a$ is
    disjoint from region $b$, then $a$ is disjoint from every region
    properly inside $b$.  We show that this is the \emph{only}
    self-absorption failure among the four non-$\EQ$ relations.

  \item \textbf{Tableau-internal resolution}
    (Section~\ref{sec:tableau-resolution}).  We observe that the
    composition filtering condition~(CF) in the tableau
    of~\cite{CompanionPaper} \emph{already forces} the edge
    reassignments needed to handle the self-absorption failure: when
    node~$n$ has a $\DR$-witness~$w$ and a $\PPI$-neighbor~$m$,
    (CF)~forces $\mathcal{E}(w, m) = \DR$, and the $\forall$-rule
    propagates the corresponding $\forall\DR$-consequences.  This
    ``enrichment'' is automatic in the original tableau---no
    modification is needed.

  \item \textbf{Direct model construction}
    (Section~\ref{sec:model-construction}).  We construct a model
    directly from the tree unraveling of an open completion graph,
    assigning cross-subtree edges via a disjunctive constraint network
    whose path-consistency follows from the tableau's properties.  The
    full RCC5 tractability theorem then yields a globally consistent
    atomic refinement, and the $\forall$-safety of all edges is
    guaranteed by the automatic enrichment.

  \item \textbf{Decidability} (Section~\ref{sec:decidability}).
    Combining our soundness result with the termination and
    completeness theorems from~\cite{CompanionPaper}, we conclude that
    $\ALCIRCC{5}$ concept satisfiability is decidable in \EXPTIME.
\end{enumerate}


%% ============================================================
\section{Preliminaries}
\label{sec:prelim}

We assume familiarity with the companion paper~\cite{CompanionPaper}
and recall only the definitions essential for the present argument.

\subsection{$\ALCIRCC{5}$ in Brief}

The logic $\ALCIRCC{5}$ has role names $\NR = \{\DR, \PO, \EQ, \PP,
\PPI\}$ with the RCC5 composition table (Table~\ref{tab:rcc5}).
Concepts are built from concept names using $\sqcap$, $\sqcup$,
$\neg$, $\exists R.C$, and $\forall R.C$ (in negation normal form;
inverse roles absorbed: $\exists \PP^{-}.C = \exists \PPI.C$).

An interpretation $\interp = (\DeltaI, \cdot^{\interp})$ is a
complete graph: every pair of distinct elements is related by exactly
one base relation from $\NR \setminus \{\EQ\}$, with $\EQ$ reserved
for identity (\emph{strong EQ} semantics).  The composition table is
enforced: $R^{\interp} \circ S^{\interp} \subseteq \bigcup_{T \in
\comp(R,S)} T^{\interp}$.

After EQ-normalization~\cite[Lemma~2.3]{CompanionPaper}, the closure
$\cl(C_0)$ contains no EQ-modalities, and we work with the four
non-EQ relations $\NR^{-} = \{\DR, \PO, \PP, \PPI\}$.

\begin{table}[t]
\centering
\caption{The RCC5 composition table.  Entry for row~$S$, column~$R$
gives $\comp(R, S)$.  The symbol $*$ denotes all five relations.}
\label{tab:rcc5}
\smallskip
\renewcommand{\arraystretch}{1.15}
\begin{tabular}{c|ccccc}
\toprule
$\comp$ & $\DR(a,b)$ & $\PO(a,b)$ & $\EQ(a,b)$ & $\PPI(a,b)$ & $\PP(a,b)$ \\
\midrule
$\DR(b,c)$  & $*$              & $\DR,\PO,\PPI$   & $\DR$  & $\DR,\PO,\PPI$       & $\DR$ \\
$\PO(b,c)$  & $\DR,\PO,\PP$    & $*$               & $\PO$  & $\PO,\PPI$           & $\DR,\PO,\PP$ \\
$\EQ(b,c)$  & $\DR$            & $\PO$             & $\EQ$  & $\PPI$               & $\PP$ \\
$\PP(b,c)$  & $\DR,\PO,\PP$    & $\PO,\PP$         & $\PP$  & $\PO,\EQ,\PP,\PPI$   & $\PP$ \\
$\PPI(b,c)$ & $\DR$            & $\DR,\PO,\PPI$    & $\PPI$ & $\PPI$               & $*$ \\
\bottomrule
\end{tabular}
\end{table}

\subsection{The Patchwork Property and Full Tractability}

\begin{theorem}[Renz \& Nebel~\cite{RenzNebel1999}; Renz~\cite{Renz1999}]
\label{thm:patchwork}
For RCC5:
\begin{enumerate}[nosep,label=(\roman*)]
  \item An atomic constraint network is consistent iff it is
    path-consistent.
  \item \textup{(}Full tractability\textup{)} A disjunctive constraint
    network is consistent iff it is path-consistent.
\end{enumerate}
\end{theorem}

\subsection{The Tableau Calculus}

We use the tableau calculus from~\cite[Section~7]{CompanionPaper}.
A \emph{completion graph} $\mathcal{G} = (V, \mathcal{L},
\mathcal{E}, \prec)$ is a complete graph with node labels
$\mathcal{L}(x) \subseteq \cl(C_0)$ and edge labels $\mathcal{E}(x,y)
\in \NR^{-}$ for distinct $x, y$.  The expansion rules are:
\begin{itemize}[nosep]
  \item $(\sqcap)$, $(\sqcup)$: propositional decomposition.
  \item $(\forall)$: if $\forall R.D \in \mathcal{L}(x)$ and
    $\mathcal{E}(x,y) = R$, add $D$ to $\mathcal{L}(y)$.
  \item $(\exists)$: if $\exists R.D \in \mathcal{L}(x)$, $x$ is
    active, and no witness exists, create a fresh node~$y$ with
    $\mathcal{L}(y) = \{D\}$, $\mathcal{E}(x,y) = R$, and
    nondeterministically assign $\mathcal{E}(z,y)$ for all other~$z$.
\end{itemize}
\emph{Blocking}: node~$x$ is blocked if some non-blocked $y \prec x$
has $\mathcal{L}(y) = \mathcal{L}(x)$.

\emph{Composition filtering} (CF): for every triple of distinct nodes
$(u, v, w)$,
$\mathcal{E}(u,w) \in \comp(\mathcal{E}(u,v), \mathcal{E}(v,w))$.

A completion graph is \emph{open} if it is saturated (no rule
applies) and clash-free (no $\{A, \neg A\} \subseteq \mathcal{L}(x)$
and (CF) holds everywhere).

\begin{theorem}[\cite{CompanionPaper}]
\label{thm:term-complete}
The tableau terminates in $2^{O(|C_0|)}$ steps.  If $C_0$ is
satisfiable, some sequence of nondeterministic choices yields an open
completion graph \textup{(}completeness\textup{)}.
\end{theorem}

\subsection{The Extension Gap}

The companion paper's soundness proof proceeds indirectly: extract
a quasimodel from the open completion graph, then build a model via
a Henkin construction that adds elements one at a time, solving a
disjunctive RCC5 CSP at each step.  The gap is that the CSP may not
be path-consistent for abstract quasimodels, causing the construction
to fail.  The present paper bypasses this indirect route.


%% ============================================================
\section{Root Cause: The Self-Absorption Failure}
\label{sec:root-cause}

The extension gap arises from a single algebraic property of the RCC5
composition table.

\begin{definition}[Self-absorption]
\label{def:self-absorption}
A base relation $R \in \NR^{-}$ is \emph{self-absorbing under}~$S \in
\NR^{-}$ if $R \in \comp(S, R)$.  Equivalently: if $S(a,b)$ and
$R(b,c)$, then $R(a,c)$ is a possible outcome.  When this fails,
knowing that $a$ is $S$-related to~$b$ \emph{excludes} the
possibility that $a$ is $R$-related to~$c$, even though $R(b,c)$
holds.
\end{definition}

\begin{lemma}[Unique self-absorption failure --- \textbf{corrected}]
\label{lem:unique-failure}
Among all pairs $(S, R) \in \NR^{-} \times \NR^{-}$, self-absorption
$R \in \comp(S, R)$ fails for exactly one pair:
\[
  \PPI \notin \comp(\DR, \PPI) = \{\DR\}.
\]
The dual under the inverse map is $\PP \notin \comp(\PP, \DR) =
\{\DR\}$ (note: $\comp(\PP, \DR)$, not $\comp(\DR, \PP)$).  For all
other $(S, R)$ pairs, $R \in \comp(S,R)$.

\smallskip\noindent
\textbf{Erratum:} The original version stated
$\PP \notin \comp(\DR, \PP) = \{\DR\}$, but the correct value is
$\comp(\DR, \PP) = \{\DR, \PO, \PP\}$ (and $\PP \in \comp(\DR, \PP)$).
The paper confused $\comp(\DR, \PP)$ with $\comp(\PP, \DR)$.
\end{lemma}

\begin{proof}
By exhaustive inspection of the composition table
(Table~\ref{tab:rcc5}).  For each $S \in \NR^{-}$ and $R \in
\NR^{-}$, we check whether $R \in \comp(S, R)$:

\smallskip
\begin{center}
\renewcommand{\arraystretch}{1.1}
\begin{tabular}{c|cccc}
\toprule
$S \setminus R$ & $\DR$ & $\PO$ & $\PP$ & $\PPI$ \\
\midrule
$\DR$  & $\checkmark$ & $\checkmark$ & $\times$ & $\times$ \\
$\PO$  & $\checkmark$ & $\checkmark$ & $\checkmark$ & $\checkmark$ \\
$\PP$  & $\checkmark$ & $\checkmark$ & $\checkmark$ & $\checkmark$ \\
$\PPI$ & $\checkmark$ & $\checkmark$ & $\checkmark$ & $\checkmark$ \\
\bottomrule
\end{tabular}
\end{center}

\smallskip\noindent
\textbf{Erratum:} The original proof claimed two failures in the
$S = \DR$ row.  The failure $\comp(\DR, \PPI) = \{\DR\}$
(row~$\PPI$, column~$\DR$) is correct.  But the claimed failure
$\comp(\DR, \PP) = \{\DR\}$ is \textbf{wrong}: the correct value is
$\comp(\DR, \PP) = \{\DR, \PO, \PP\}$ (row~$\PP$, column~$\DR$),
and $\PP \in \comp(\DR, \PP)$.  The correct dual is
$\PP \notin \comp(\PP, \DR) = \{\DR\}$.
\end{proof}

\begin{remark}[Geometric interpretation]
\label{rem:geometric}
The identity $\comp(\DR, \PPI) = \{\DR\}$ has a transparent geometric
meaning: if region~$a$ is \emph{disjoint} from region~$b$, and
region~$c$ is \emph{properly inside}~$b$ (i.e., $\PPI(b,c)$), then
$a$ must be disjoint from~$c$ as well---$a$ cannot ``reach inside''
$b$ to touch~$c$.  We call this the \emph{containment collapse}:
the $\DR$ relation propagates inward through $\PPI$-containment.

The dual $\comp(\PP, \DR) = \{\DR\}$ says the same from the inside
out: if $a$ is properly inside $b$ ($\PP(a,b)$) and $b$ is disjoint
from $c$ ($\DR(b,c)$), then $a$ is disjoint from~$c$.
\end{remark}

\begin{lemma}[Containment collapse]
\label{lem:containment-collapse}
For all $k \geq 1$:
\[
  \comp(\DR, \underbrace{\PPI, \ldots, \PPI}_{k}) = \{\DR\}.
\]
That is, $\DR$ composed with any number of $\PPI$ steps yields only
$\DR$.
\end{lemma}

\begin{proof}
By induction.  Base case: $\comp(\DR, \PPI) = \{\DR\}$.  Inductive
step: $\comp(\{\DR\}, \PPI) = \comp(\DR, \PPI) = \{\DR\}$.
\end{proof}

\subsection{Why Self-Absorption Matters for Model Construction}

The self-absorption failure is the root cause of the extension gap.
Consider the model construction from an open completion graph.  The
standard approach ``unravels'' blocked nodes: when node~$x$ is
blocked by~$y$ (same type), $x$'s subtree is replaced by a copy of
$y$'s subtree, creating infinitely many domain elements.  We call the
resulting infinite structure the \emph{tree unraveling} of the
completion graph.

Each domain element~$d$ in the unraveling maps to a tableau node
$T(d)$, inheriting its type $\mathcal{L}(T(d))$.  For the unraveling
to yield a valid model, we must assign edges $R(d_1, d_2)$ between
\emph{all} pairs of domain elements (not just tree-adjacent ones),
satisfying:
\begin{enumerate}[nosep,label=(M\arabic*)]
  \item \textbf{Composition consistency}: for all triples $(d_1, d_2,
    d_3)$, $R(d_1, d_3) \in \comp(R(d_1, d_2), R(d_2, d_3))$.
  \item \textbf{$\forall$-safety}: for every edge $R(d_1, d_2) = S$,
    if $\forall S.D \in \mathcal{L}(T(d_1))$ then $D \in
    \mathcal{L}(T(d_2))$.
\end{enumerate}

The natural strategy is \emph{ancestor projection}: for $d_1, d_2$
with $T(d_1) \neq T(d_2)$, set $R(d_1, d_2) = \mathcal{E}(T(d_1),
T(d_2))$.  This automatically satisfies (M2) (the tableau's
$(\forall)$-rule ensures $\forall$-safety for all tableau edges) and
satisfies (M1) for triples mapping to three distinct tableau nodes
(by (CF) in the completion graph).

The problem arises for triples $(d_1, d_2, d_3)$ where $T(d_1) =
T(d_2) = n$ (two copies of the same tableau node) and $T(d_3) = m
\neq n$.  Ancestor projection gives:
\[
  R(d_1, d_3) = R(d_2, d_3) = \mathcal{E}(n, m).
\]
If $R(d_1, d_2) = S$ (the self-safe relation between copies),
composition consistency requires:
\begin{equation}
\label{eq:self-absorption}
  \mathcal{E}(n, m) \in \comp(S, \mathcal{E}(n, m)).
\end{equation}
This is precisely the self-absorption condition.  When $S = \DR$ and
$\mathcal{E}(n, m) = \PPI$, equation~\eqref{eq:self-absorption}
fails: $\PPI \notin \comp(\DR, \PPI) = \{\DR\}$.


%% ============================================================
\section{Tableau-Internal Resolution}
\label{sec:tableau-resolution}

We now show that the tableau's own composition filtering (CF) already
contains the information needed to resolve the self-absorption failure.
The key observation is that the problematic configuration---a $\DR$-witness
alongside a $\PPI$-neighbor---forces a specific edge assignment within the
completion graph, and the $(\forall)$-rule propagates the consequences.

\subsection{Forced Edge Reassignment}

\begin{lemma}[Forced $\DR$ edge]
\label{lem:forced-DR}
Let $\mathcal{G}$ be a completion graph satisfying~\textup{(CF)}.
Suppose nodes $n, w, m$ satisfy $\mathcal{E}(n, w) = \DR$ and
$\mathcal{E}(n, m) = \PPI$.  Then $\mathcal{E}(w, m) = \DR$.
\end{lemma}

\begin{proof}
By (CF) on the triple $(n, w, m)$:
\[
  \mathcal{E}(n, m) \in \comp(\mathcal{E}(n, w),\, \mathcal{E}(w, m))
  = \comp(\DR,\, \mathcal{E}(w, m)).
\]
We need $\PPI \in \comp(\DR, \mathcal{E}(w, m))$.  Checking each
possible value of $\mathcal{E}(w, m)$:
\begin{align*}
  \comp(\DR, \DR) &= \NR^{-} \cup \{\EQ\}, \quad \PPI \in \checkmark \\
  \comp(\DR, \PO) &= \{\DR, \PO, \PP\}, \quad \PPI \notin \text{\;---} \\
  \comp(\DR, \PP) &= \{\DR, \PO, \PP\}, \quad \PPI \notin \text{\;---} \\
  \comp(\DR, \PPI) &= \{\DR\}, \quad \PPI \notin \text{\;---}
\end{align*}
Only $\mathcal{E}(w, m) = \DR$ satisfies the constraint.
\end{proof}

\begin{corollary}[Automatic enrichment]
\label{cor:enrichment}
In the setting of Lemma~\ref{lem:forced-DR}, the $(\forall)$-rule
applied to the edge $\mathcal{E}(w, m) = \DR$ ensures:
\begin{enumerate}[nosep,label=(\alph*)]
  \item $\forall \DR.D \in \mathcal{L}(w) \;\Rightarrow\; D \in
    \mathcal{L}(m)$,
  \item $\forall \DR.D \in \mathcal{L}(m) \;\Rightarrow\; D \in
    \mathcal{L}(w)$.
\end{enumerate}
These are the ``enrichment'' conditions that the model construction
requires for $\DR$-edges between copies of~$n$ and elements
corresponding to~$m$.  They are satisfied \emph{automatically} by the
saturated completion graph, without any modification to the tableau.
\end{corollary}

\begin{remark}[Dual case]
\label{rem:dual}
The dual scenario $\mathcal{E}(n, w) = \DR$ and $\mathcal{E}(n, m)
= \PP$ forces $\mathcal{E}(w, m) \in \comp(\DR, \mathcal{E}(w, m))$
with $\PP \in \comp(\DR, \mathcal{E}(w, m))$.  By a similar check,
$\mathcal{E}(w, m) = \DR$ is the only possibility, and the
$(\forall)$-rule handles it identically.  More generally,
Lemma~\ref{lem:containment-collapse} shows that $\comp(\DR,
\PPI)^k = \{\DR\}$ for all $k$: every node reachable from~$m$ by a
chain of $\PPI$-edges is forced to $\DR$ with~$w$.
\end{remark}

\subsection{Self-Safety of Witness Relations}

When two copies of the same tableau node need an edge between them, we
use the \emph{witness relation}---the role~$S$ from the existential
demand $\exists S.D$ that created the blocked node.

\begin{definition}[Self-safe relation]
\label{def:self-safe}
A relation $S \in \NR^{-}$ is \emph{self-safe} for a type $\tau
\subseteq \cl(C_0)$ if:
\begin{enumerate}[nosep,label=(\roman*)]
  \item $\forall S.D \in \tau \;\Rightarrow\; D \in \tau$, and
  \item $\forall \inv(S).D \in \tau \;\Rightarrow\; D \in \tau$.
\end{enumerate}
Equivalently: using~$S$ as an edge between two elements of type~$\tau$
does not violate any $\forall$-constraint.
\end{definition}

\begin{theorem}[Witness relations are self-safe]
\label{thm:self-safe}
Let $\mathcal{G}$ be a saturated completion graph.  Suppose active
node~$n$ has demand $\exists S.D$, witnessed by node~$w$ with
$\mathcal{E}(n, w) = S$ and $D \in \mathcal{L}(w)$.  If $w$ is
blocked by~$n$ \textup{(}i.e., $\mathcal{L}(w) = \mathcal{L}(n) =
\tau$\textup{)}, then $S$ is self-safe for~$\tau$.
\end{theorem}

\begin{proof}
Since the graph is saturated, the $(\forall)$-rule has been
exhaustively applied:
\begin{itemize}[nosep]
  \item Edge $\mathcal{E}(n, w) = S$: for every $\forall S.D' \in
    \mathcal{L}(n) = \tau$, the $(\forall)$-rule gives $D' \in
    \mathcal{L}(w) = \tau$.  This establishes condition~(i).
  \item Edge $\mathcal{E}(w, n) = \inv(S)$: for every $\forall
    \inv(S).D' \in \mathcal{L}(w) = \tau$, the $(\forall)$-rule
    gives $D' \in \mathcal{L}(n) = \tau$.  This establishes
    condition~(ii). \qedhere
\end{itemize}
\end{proof}

\begin{lemma}[Iterated self-composition]
\label{lem:iterated-comp}
For every $S \in \NR^{-}$: $S \in \comp(S, S)$.  Consequently,
$S \in \comp(S)^k$ for all $k \geq 1$ \textup{(}by induction\textup{)}.
\end{lemma}

\begin{proof}
By inspection of the composition table:
\begin{align*}
  \DR &\in \comp(\DR, \DR) = \NR^{-} \cup \{\EQ\}, \quad\checkmark \\
  \PO &\in \comp(\PO, \PO) = \NR^{-} \cup \{\EQ\}, \quad\checkmark \\
  \PP &\in \comp(\PP, \PP) = \{\PP\}, \quad\checkmark \\
  \PPI &\in \comp(\PPI, \PPI) = \{\PPI\}, \quad\checkmark \qedhere
\end{align*}
\end{proof}

\begin{corollary}
\label{cor:copy-edge}
In the tree unraveling, if $d_1, d_2, \ldots, d_k$ are consecutive
copies of the same tableau node~$n$, each linked by witness
relation~$S$, then the edge $R(d_i, d_j) = S$ is
composition-consistent with the tree edges \textup{(}by
Lemma~\ref{lem:iterated-comp}\textup{)} and $\forall$-safe
\textup{(}by Theorem~\ref{thm:self-safe}\textup{)}.
\end{corollary}

\subsection{Universal Self-Absorption for Non-DR Witnesses}

\begin{lemma}
\label{lem:non-DR-ok}
For $S \in \{\PO, \PP, \PPI\}$: $R \in \comp(S, R)$ for all $R \in
\NR^{-}$.
\end{lemma}

\begin{proof}
From Lemma~\ref{lem:unique-failure}: the only self-absorption
failures are $(S, R) = (\DR, \PP)$ and $(S, R) = (\DR, \PPI)$.  For
$S \neq \DR$, there are no failures.
\end{proof}

This means that for witness relations $S \in \{\PO, \PP, \PPI\}$,
\emph{pure ancestor projection works without modification}.  The copy
edges use relation~$S$ between copies, and for any cross-subtree edge
$R(d_1, d_3) = \mathcal{E}(n, m)$, self-absorption gives
$\mathcal{E}(n, m) \in \comp(S, \mathcal{E}(n, m))$.  Composition
consistency of the triple $(d_1, d_2, d_3)$ is immediate.

\emph{Only DR-witnesses require the modified edge assignment described
in Lemma~\ref{lem:forced-DR}.}


%% ============================================================
\section{Direct Model Construction}
\label{sec:model-construction}

We now construct a model from an open completion graph without passing
through the quasimodel abstraction.

\subsection{The Tree Unraveling}

\begin{definition}[Tree unraveling]
\label{def:unraveling}
Let $\mathcal{G} = (V, \mathcal{L}, \mathcal{E}, \prec)$ be an open
completion graph for~$C_0$.  Let $V_a \subseteq V$ be the set of
active (non-blocked) nodes.  For each blocked node~$x$, let
$\beta(x) \in V_a$ denote its blocker ($\mathcal{L}(\beta(x)) =
\mathcal{L}(x)$, $\beta(x) \prec x$, $\beta(x)$ not blocked).

The \emph{tree unraveling} $\mathfrak{T} = (\Delta, T, \mathsf{par},
\sigma)$ is defined inductively:
\begin{enumerate}[nosep]
  \item The root element $d_0 \in \Delta$ has $T(d_0) = x_0$ (the
    initial tableau node).
  \item For each element $d \in \Delta$ with $T(d) = n$, and each
    unsatisfied demand $\exists S.D \in \mathcal{L}(n)$ witnessed by
    node~$w$ (i.e., $\mathcal{E}(n, w) = S$ and $D \in
    \mathcal{L}(w)$): create a child element~$d'$ with
    $\mathsf{par}(d') = d$, $\sigma(d') = S$
    (the \emph{witness relation}), and:
    \begin{itemize}[nosep]
      \item If $w$ is active: $T(d') = w$.
      \item If $w$ is blocked by $y = \beta(w)$: $T(d') = y$.
    \end{itemize}
\end{enumerate}
The map $T \colon \Delta \to V_a$ assigns each domain element its
\emph{corresponding active node}.  The function $\sigma(d)$ records
the witness relation used to create~$d$.
\end{definition}

\begin{remark}
Since the completion graph is saturated and every active node's
demands are witnessed, the unraveling is well-defined.  Blocking
creates ``loops'' that make $\mathfrak{T}$ an infinite tree whenever
the completion graph has blocked nodes.  Each element~$d$ inherits its
type: $\mathsf{tp}(d) = \mathcal{L}(T(d))$.
\end{remark}

\subsection{Edge Assignment}

We define a model $\interp = (\Delta, \cdot^{\interp})$ on the domain
$\Delta$ of the tree unraveling.

\begin{definition}[Edge assignment]
\label{def:edge-assignment}
For each pair of distinct elements $d_1, d_2 \in \Delta$, we define
$R^{\interp}(d_1, d_2) \in \NR^{-}$ as follows.

\medskip
\noindent\textbf{Case 1: Parent--child pair.}\;
If $\mathsf{par}(d_2) = d_1$ (or vice versa), set $R^{\interp}(d_1,
d_2) = \sigma(d_2)$ (the witness relation).

\medskip
\noindent\textbf{Case 2: Non-adjacent pair.}\;
For all other pairs, we construct a disjunctive constraint network
$\mathcal{N}$ and invoke Theorem~\ref{thm:patchwork}(ii) to obtain a
consistent atomic refinement.

The network $\mathcal{N}$ assigns to each non-adjacent pair $(d_1,
d_2)$ the domain:
\begin{equation}
\label{eq:domain}
  D(d_1, d_2) = \DN_{\mathsf{safe}}(\mathsf{tp}(d_1),\,
    \mathsf{tp}(d_2))
\end{equation}
where
\[
  \DN_{\mathsf{safe}}(\tau_1, \tau_2) \;=\; \bigl\{
    \mathcal{E}(a, b) : a, b \in V,\;
    \mathcal{L}(a) = \tau_1,\; \mathcal{L}(b) = \tau_2,\; a \neq b
  \bigr\}.
\]
This is the set of all base relations that \emph{actually appear} as
edges between nodes of types $\tau_1$ and $\tau_2$ in the completion
graph.  Since the $(\forall)$-rule has been exhaustively applied, every
such relation is $\forall$-safe.
\end{definition}

\begin{remark}
The domain $\DN_{\mathsf{safe}}(\tau_1, \tau_2)$ is non-empty for all
type pairs $(\tau_1, \tau_2)$ appearing in the completion graph: if
$\tau_1 = \tau_2 = \tau$, a single active node~$a$ of type~$\tau$
does not provide a pair, but two distinct nodes of the same type
(one active, one blocked) provide an edge.  If $\tau_1 \neq \tau_2$,
the unique active nodes of each type provide an edge.

For same-type pairs where only one node of that type exists in the
completion graph, $\DN_{\mathsf{safe}}(\tau, \tau)$ may need to
include self-safe relations.  By Theorem~\ref{thm:self-safe}, the
witness relation~$S$ used for blocking is always available.
\end{remark}

\subsection{Path-Consistency of the Constraint Network}

The central technical result is that the constraint network
$\mathcal{N}$ is path-consistent.

\begin{theorem}[\textbf{RETRACTED} --- Path-consistency]
\label{thm:path-consistent}
\sout{The disjunctive constraint network $\mathcal{N}$ defined in
Definition~\ref{def:edge-assignment} \textup{(}with the parent--child
edges fixed as in Case~1\textup{)} is path-consistent.}
\end{theorem}

\noindent
\textbf{This theorem is false.}  The $\DN_{\mathsf{safe}}$ domains
aggregate relations from different representatives of the same type.
When two nodes of type $\tau_1$ have different relations to a node of
type $\tau_2$, the domain contains both, and arc-consistency
enforcement can empty domains.  GPT-5.4 provided a concrete
counterexample.  See the status box for details.

\begin{proof}
We must show: for every triple of distinct elements $(d_1, d_2, d_3)$
and every $R_{12} \in D(d_1, d_2)$, there exists $R_{23} \in D(d_2,
d_3)$ with $\comp(R_{12}, R_{23}) \cap D(d_1, d_3) \neq \emptyset$.

We analyze cases based on how many of the three pairs are
parent--child (Case~1) edges.

\medskip
\noindent\textbf{Case A: All three pairs are non-adjacent
(Case~2).}\;
The domains are $D(d_i, d_j) = \DN_{\mathsf{safe}}(\mathsf{tp}(d_i),
\mathsf{tp}(d_j))$.  Let $\tau_i = \mathsf{tp}(d_i)$ for $i = 1,2,3$.

For $R_{12} \in D(d_1, d_2)$, there exist nodes $a_1$ (type~$\tau_1$)
and $a_2$ (type~$\tau_2$) in $\mathcal{G}$ with $\mathcal{E}(a_1,
a_2) = R_{12}$.  Let $a_3$ be any node of type~$\tau_3$.  By (CF) on
$(a_1, a_2, a_3)$:
\[
  \mathcal{E}(a_1, a_3) \in \comp(\mathcal{E}(a_1, a_2),\,
    \mathcal{E}(a_2, a_3)) = \comp(R_{12},\, \mathcal{E}(a_2, a_3)).
\]
Set $R_{23} = \mathcal{E}(a_2, a_3)$ and $R_{13} = \mathcal{E}(a_1,
a_3)$.  Then $R_{23} \in D(d_2, d_3)$, $R_{13} \in D(d_1, d_3)$,
and $R_{13} \in \comp(R_{12}, R_{23})$.

\medskip
\noindent\textbf{Case B: Exactly one parent--child pair.}\;
Say $d_2 = \mathsf{par}(d_3)$ with $R_{23} = \sigma(d_3) = S$ fixed.
Then $T(d_2)$ has a witness~$w$ with $\mathcal{E}(T(d_2), w) = S$
and $T(d_3)$ is the active node corresponding to~$w$ (or its
blocker).

For $R_{12} \in D(d_1, d_2)$, there exist $a_1$ (type~$\tau_1$) and
$a_2$ (type~$\tau_2$) with $\mathcal{E}(a_1, a_2) = R_{12}$.  We
need $\comp(R_{12}, S) \cap D(d_1, d_3) \neq \emptyset$.

In the completion graph, $T(d_2)$ has edge~$S$ to some node~$w'$
of type~$\tau_3 = \mathsf{tp}(d_3)$.  By (CF) on $(a_1, T(d_2), w')$:
$\mathcal{E}(a_1, w') \in \comp(\mathcal{E}(a_1, T(d_2)), S)$.

If $a_1 = T(d_1)$ or has the same type, $\mathcal{E}(a_1, T(d_2))$
might differ from $R_{12}$.  However, by choosing $a_1$ and~$a_2$ to
be the specific pair realizing $R_{12}$, and using (CF) on $(a_1, a_2,
w')$ where $w'$ is the witness of $T(d_2)$'s demand:
\[
  \mathcal{E}(a_1, w') \in \comp(\mathcal{E}(a_1, a_2),\,
    \mathcal{E}(a_2, w')).
\]
Now, $\mathcal{E}(a_2, w')$ might not equal~$S$ (since $a_2$ is a
node of type~$\tau_2$ but may differ from $T(d_2)$).  If $a_2 =
T(d_2)$, then $\mathcal{E}(a_2, w') = S$ and we are done.  If $a_2
\neq T(d_2)$, we use the (CF) constraint on $(T(d_2), a_2, w')$:
$\mathcal{E}(T(d_2), w') = S$ and
$\mathcal{E}(a_2, w') \in \comp(\mathcal{E}(a_2, T(d_2)),
S)$.  In general, the (CF) constraints on the completion graph
ensure that for any $R_{12} \in D(d_1, d_2)$, there exists
$R_{13} \in D(d_1, d_3)$ with $R_{13} \in \comp(R_{12}, S)$.
(We prove this by noting that the pair $(a_1, w')$ provides a witness
in $D(d_1, d_3)$.)

\medskip
\noindent\textbf{Case C: Two or three parent--child pairs.}\;
When two edges are fixed tree edges, the third must satisfy the
composition constraint.  By (CF) in the completion graph applied to the
corresponding tableau nodes, the constraint is satisfiable, and the
remaining domain $D$ is non-empty.  The argument mirrors Case~B with
two fixed relations instead of one.
\end{proof}

\begin{corollary}[Consistent edge assignment exists]
\label{cor:consistent-assignment}
The disjunctive constraint network $\mathcal{N}$ has a consistent
atomic refinement: there exists a function $R^{\interp} \colon \Delta
\times \Delta \to \NR$ such that $R^{\interp}(d, d) = \EQ$, and for
distinct $d_1, d_2$, $R^{\interp}(d_1, d_2) \in D(d_1, d_2)$ and the
resulting atomic network is globally consistent.
\end{corollary}

\begin{proof}
By Theorem~\ref{thm:path-consistent} and the full RCC5 tractability
(Theorem~\ref{thm:patchwork}(ii)).
\end{proof}

\subsection{Verification of Model Properties}

\begin{theorem}[\textbf{RETRACTED} --- Soundness]
\label{thm:soundness}
\sout{If the tableau produces an open completion graph $\mathcal{G}$ for
$C_0$, then $C_0$ is satisfiable in $\ALCIRCC{5}$.}
\end{theorem}

\noindent
\textbf{Retracted:} depends on Theorem~\ref{thm:path-consistent},
which is false.

\begin{proof}
We verify that the interpretation $\interp = (\Delta,
\cdot^{\interp})$ constructed above is an $\ALCIRCC{5}$ model with
$d_0 \in C_0^{\interp}$.

\medskip
\noindent\textbf{(S1) One cluster.}\;
$R^{\interp}$ assigns a relation to every pair, by construction.

\medskip
\noindent\textbf{(S2) JEPD.}\;
Each pair gets exactly one base relation ($\EQ$ for self-pairs,
exactly one relation from $\NR^{-}$ for distinct pairs).

\medskip
\noindent\textbf{(S3) Composition.}\;
Global consistency of the atomic network (Corollary~\ref{cor:consistent-assignment})
implies: for all distinct $d_1, d_2, d_3$,
$R^{\interp}(d_1, d_3) \in \comp(R^{\interp}(d_1, d_2),
R^{\interp}(d_2, d_3))$.

\medskip
\noindent\textbf{(S4) Converse.}\;
The edge assignment respects converse: $R^{\interp}(d_2, d_1) =
\inv(R^{\interp}(d_1, d_2))$ (inherited from the disjunctive network's
converse symmetry and the completion graph's converse constraint).

\medskip
\noindent\textbf{(S5) Strong EQ.}\;
$R^{\interp}(d, d) = \EQ$ for all~$d$, and distinct elements get
non-$\EQ$ relations.

\medskip
\noindent\textbf{$\forall$-safety.}\;
For every edge $R^{\interp}(d_1, d_2) = S$:
since $S \in D(d_1, d_2) = \DN_{\mathsf{safe}}(\mathsf{tp}(d_1),
\mathsf{tp}(d_2))$, there exist nodes $a_1, a_2$ in $\mathcal{G}$
with $\mathcal{L}(a_1) = \mathsf{tp}(d_1)$, $\mathcal{L}(a_2) =
\mathsf{tp}(d_2)$, and $\mathcal{E}(a_1, a_2) = S$.  The
$(\forall)$-rule ensures: $\forall S.D \in \mathcal{L}(a_1)
\Rightarrow D \in \mathcal{L}(a_2)$.  Since $\mathcal{L}(a_i) =
\mathsf{tp}(d_i)$, this gives $\forall S.D \in \mathsf{tp}(d_1)
\Rightarrow D \in \mathsf{tp}(d_2)$.

\medskip
\noindent\textbf{$\exists$-satisfaction.}\;
For every $d \in \Delta$ and $\exists S.D \in \mathsf{tp}(d)$: by
construction of the tree unraveling, $d$ has a child~$d'$ with
$\sigma(d') = S$ and $D \in \mathsf{tp}(d')$ (the demand was
witnessed in the saturated completion graph).  The tree edge gives
$R^{\interp}(d, d') = S$, so $d'$ is an $S$-neighbor of~$d$ with
$D \in \mathsf{tp}(d')$.  Since all formulas in $\mathsf{tp}(d')$
hold in the model (by induction on the closure, using the
$\forall$-safety verified above), $d' \in D^{\interp}$.

\medskip
\noindent\textbf{Propositional correctness.}\;
Types are inherited from the completion graph, where propositional
rules $(\sqcap)$, $(\sqcup)$ were exhaustively applied.  Since
$\mathsf{tp}(d)$ is a subset of $\cl(C_0)$ that is
propositionally closed and clash-free, standard induction on concept
structure shows $d \in E^{\interp}$ for every $E \in \mathsf{tp}(d)$.

\medskip
\noindent\textbf{Root element.}\;
$\mathsf{tp}(d_0) = \mathcal{L}(x_0) \ni C_0$, so $d_0 \in
C_0^{\interp}$.
\end{proof}


%% ============================================================
\section{Decidability of $\ALCIRCC{5}$}
\label{sec:decidability}

\begin{theorem}[\textbf{RETRACTED} --- Main theorem]
\label{thm:main}
\sout{Concept satisfiability in $\ALCIRCC{5}$ is decidable.  The tableau
calculus of~\cite{CompanionPaper} is a sound and complete decision
procedure running in \EXPTIME.}
\end{theorem}

\noindent
\textbf{Retracted:} depends on Theorem~\ref{thm:soundness}, which
depends on Theorem~\ref{thm:path-consistent}, which is false.
Additionally, the companion paper's type-elimination algorithm is
unsound (rejects satisfiable concepts), and the companion paper's
tableau soundness itself depends on the Henkin construction, which
has the extension gap that this paper attempted to close.

\begin{corollary}[\textbf{RETRACTED}]
\label{cor:complexity}
\sout{$\PSPACE \leq \ALCIRCC{5} \leq \EXPTIME$.}
\end{corollary}

\noindent
\textbf{Retracted:} the \PSPACE{} lower bound
(Wessel~\cite{Wessel2003}) remains valid, but the \EXPTIME{} upper
bound is not established.


%% ============================================================
\section{Discussion}
\label{sec:discussion}

\subsection{What Was the Extension Gap?}

The extension gap in~\cite{CompanionPaper} arose from the
\emph{indirect} soundness proof: open completion graph $\to$
quasimodel $\to$ model (via Henkin construction).  The quasimodel
abstraction discards the specific node-to-node edges of the completion
graph, retaining only the \emph{set} of pair-types between type pairs.
This loss of information made it impossible to verify that the Henkin
construction's CSP is path-consistent at each step.

Our resolution bypasses the quasimodel abstraction entirely.  The
direct model construction from the tree unraveling retains the full
edge information of the completion graph, which provides enough
structure to establish path-consistency of the disjunctive constraint
network.

\subsection{The Role of the Self-Absorption Failure}

The unique self-absorption failure $\PPI \notin \comp(\DR, \PPI) =
\{\DR\}$ is the algebraic root cause of the extension gap, but it is
\emph{not} an obstacle to decidability.  Rather, it is automatically
resolved by the composition filtering condition~(CF) within the
tableau: the forced $\DR$-edge (Lemma~\ref{lem:forced-DR}) and the
resulting $\forall$-rule propagation (Corollary~\ref{cor:enrichment})
ensure that the types of relevant nodes already contain all the
$\forall\DR$-consequences needed for the model construction.

The geometric content of this resolution is natural: if a region~$a$
is disjoint from region~$b$ (the $\DR$-witness), and region~$m$ is
inside~$b$ (the $\PPI$-neighbor), then the model construction
acknowledges that $a$ is necessarily disjoint from~$m$ and ensures
the concept labels reflect this.

\subsection{Extension to RCC8}

The RCC8 composition table has its own self-absorption failures (e.g.,
involving $\mathsf{DC}$, $\mathsf{NTPP}$, and $\mathsf{NTPPI}$).
The forced-edge argument (Lemma~\ref{lem:forced-DR}) generalizes: the
tableau's (CF) forces specific edges in the problematic
configurations.  However, RCC8 lacks full tractability
(Theorem~\ref{thm:patchwork}(ii) holds only for RCC5), so the
path-consistency argument in Theorem~\ref{thm:path-consistent} does
not directly transfer.  The decidability of $\ALCIRCC{8}$ may require
a different approach to the constraint solving step.

\subsection{Post-Retraction Assessment}

Theorem~\ref{thm:path-consistent} (path-consistency of the constraint
network) is \textbf{false}.  The flaw is in the $\DN_{\mathsf{safe}}$
domain definition: it aggregates all relations appearing between any
pair of nodes of given types, but different representatives of the same
type may have different relations to a third type.  This means the
completion graph's (CF) witnesses used in the proof may come from
different representatives, and the composition constraints between
these witnesses need not be jointly satisfiable.

\emph{What survives:}
\begin{itemize}[nosep]
  \item The root cause analysis (Section~\ref{sec:root-cause}):
    the unique self-absorption failure
    $\PPI \notin \comp(\DR, \PPI) = \{\DR\}$ after fixing the
    algebraic error.
  \item Lemma~\ref{lem:forced-DR} (forced $\DR$ edge) and
    Corollary~\ref{cor:enrichment} (automatic enrichment).
  \item Theorem~\ref{thm:self-safe} (witness relations are self-safe)
    and Lemma~\ref{lem:iterated-comp} (iterated self-composition).
  \item The tree unraveling construction
    (Definition~\ref{def:unraveling}).
\end{itemize}

\emph{What fails:}
\begin{itemize}[nosep]
  \item Theorem~\ref{thm:path-consistent}: $\DN_{\mathsf{safe}}$
    domains are too coarse for path-consistency.
  \item All results depending on it: Theorems~\ref{thm:soundness}
    and~\ref{thm:main}, Corollary~\ref{cor:complexity}.
\end{itemize}


%% ============================================================
\section*{Acknowledgments}

This research was prompted by Michael Wessel, who introduced the
$\ALCIRCC{}$ family during his doctoral work at the University of
Hamburg and left the decidability of $\ALCIRCC{5}$ as an open problem.
The root cause analysis (Section~\ref{sec:root-cause}) emerged from a
systematic computational investigation of the RCC5 composition table's
self-absorption properties, confirming the unique failure
$\PPI \notin \comp(\DR, \PPI)$ as the source of the extension gap.

%% ============================================================
\bibliographystyle{plain}
\begin{thebibliography}{10}

\bibitem{CompanionPaper}
Claude.
\newblock On the decidability of {$\mathcal{ALCI}_{\mathit{RCC5}}$} and
  {$\mathcal{ALCI}_{\mathit{RCC8}}$}: Quasimodels meet the patchwork
  property.
\newblock Companion paper, March 2026 (revised).

\bibitem{BaaderEtAl2003}
F.~Baader, D.~Calvanese, D.~McGuinness, D.~Nardi, and P.~Patel-Schneider,
  editors.
\newblock {\em The Description Logic Handbook: Theory, Implementation, and
  Applications}.
\newblock Cambridge University Press, 2003.

\bibitem{Renz1999}
J.~Renz.
\newblock Maximal tractable fragments of the {Region Connection Calculus}: A
  complete analysis.
\newblock In {\em Proc.\ IJCAI}, pages 448--454, 1999.

\bibitem{RenzNebel1999}
J.~Renz and B.~Nebel.
\newblock On the complexity of qualitative spatial reasoning: A maximal
  tractable fragment of the {Region Connection Calculus}.
\newblock {\em Artificial Intelligence}, 108(1-2):69--123, 1999.

\bibitem{Wessel2000}
M.~Wessel.
\newblock Decidable and undecidable extensions of {$\mathcal{ALC}$} with
  composition-based role inclusion axioms.
\newblock Technical Report FBI-HH-M-301/01, University of Hamburg, 2000.

\bibitem{Wessel2003}
M.~Wessel.
\newblock Qualitative spatial reasoning with the
  {$\mathcal{ALCI}_\mathit{RCC}$} family---{F}irst results and unanswered
  questions.
\newblock Technical Report FBI-HH-M-324/03, University of Hamburg, 2002/2003.

\end{thebibliography}

\end{document}
